% ---------------------------------------------------------------------
\section{Talep işleme hattı ve tam bilgili optimumlar}\label{sec:pipeline}
% ---------------------------------------------------------------------

\subsection{İstasyon düzeyinde talep işleme hattı}\label{app:admission}

Varış sınırlaması ve kabul kontrolü denetleyiciye değil, istasyon işletmecisinin talep ayarına aittir.
Her test tekrarı için bunlar herhangi bir denetleyici çalışmadan önce bir kez hesaplanır ve ortaya çıkan kabul edilen talepler $\Ereq_i$, karşılaştırılan her yöntem ve her iki optimum tarafından paylaşılır.
Sarmalayıcı, kabul edilen gereksinimi varışta EV2Gym'in EV nesnesine yazar; böylece EV-GNN ve Gurobi eMPC gibi harici karşılaştırma yöntemleri de aynı hedefi izler.
HiGHS eMPC, ufku içindeki varışlar için paylaşılan girdileri kullanır.
Konformal geçişler, paylaşılan gereksinimi kendi $\hat d_i^{\sm{lo}}$ sınırlarıyla yeniden sertifikalar; dolayısıyla bu geçişlerin gereksinimi hiçbir zaman daha büyük değildir; kabul kontrolü içermeyen geçişler ise varışta sınırlandırılmış talebi korur.
Her geçişte eksik enerji, paylaşılan kabul edilen gereksinime göre ölçülür.

Varışlar zaman sırasına göre işlenir; eşitlik hâlinde şarj ünitesi indisi belirleyicidir.
Varan bir EV $j$; dışsal varış enerjisi, istenen enerjisi, ayrılış adımı ve EV parametreleriyle girer ve \Eq{M-eq:cap} ile verilen sınırlama, sertifikalı bir eşik tablosuyla birlikte ilk gereksinimini verir.
Ardından kabul kontrolü, varış adımı $t_0$'da, daha önce kabul edilmiş ve hâlâ bağlı olan EV'ler ile EV $j$ üzerinde \Eq{eq:admission} ile verilen olurluluk LP'sini çözer:
\begin{equation}\label{eq:admission}
\begin{aligned}
&E_i(t+1)=E_i(t)+\Delta t\,p_i(t),\qquad E_i^{\min}\le E_i(t+1)\le E_i^{\sm{cap}},\qquad E_i(d_i+1)\ge E_i^{\sm{tgt}},\\
&-\Pdis_i\le p_i(t)\le\Pnom_i,\qquad p_i(t)\le\varrho_i^{\sm{in}}\,(E_i^{\sm{cap}}-E_i(t))\ (t>t_0),\qquad p_i(t_0)\in[\pmn_i,\bar P_i],\\
&-\hat C_{\sm{rev}}(t)\le\textstyle\sum_ip_i(t)\le\hat C(t)-n(t_0)\varepsilon_q .
\end{aligned}
\end{equation}
Önceki EV'lerin $t_0$'daki enerjileri herhangi bir denetleyicinin yörüngesinden değil, nominal bir plandan okunur.
Her kabul kararından sonra aynı LP'nin esnek bir sürümü bu planı, hedeflere olabildiğince erken ulaşılacak şekilde günceller.
$\hat C$ ve $\hat C_{\sm{rev}}$ limitleri $t_0$'daki ölçülen değerleri, sonraki $H_f$ adım için gürültülü tahminleri ve sonrasında son tahmini kullanır.
Eğim $\varrho_i^{\sm{in}}$, tüm $E$ değerleri için $\bar P_i(E)\ge\min\{\Pnom_i,\varrho_i^{\sm{in}}(E_i^{\sm{cap}}-E)\}$ koşulunu sağlayan en büyük değerdir; dolayısıyla doğrusal şarj gücü azalması modeli bir iç yaklaşıklıktır ve LP açısından olurlu her çizelge fiziksel olarak gerçeklenebilir.

LP içinde yeni EV'nin hedefi $E_j^{\sm{tgt}}=\Ereq_j$'dir; önceki bir EV'nin hedefi ise $\Ereq_i$ ile $\Psi_i^{\sm{LP}}(E_i(t_0))-10^{-6}$ değerlerinin küçüğüdür.
Burada $\Psi_i^{\sm{LP}}$, LP'nin şarj gücü azalması modelinde erişilebilen en büyük ayrılış enerjisidir ve bu hedef yalnızca kabul LP'si içinde kullanılır.
Tam batarya talepleri doğrudan hedef olarak alınsaydı, böyle bir EV bağlı olduğu her durumda LP olursuz olurdu; çünkü doğrusal şarj gücü azalması modeli sonlu sayıda adımda hiçbir zaman tam şarja ulaşmaz.
Olurluluk yeni talepte monotondur; dolayısıyla $E_j(d_j+1)$'i en büyükleyen tek bir LP, kabul edilebilir en büyük talebi verir.
Daha önce kabul edilen talepler hiçbir zaman azaltılmaz ve EV $j$'nin düşürülmüş talebi yeniden sınırlama rutininden geçer; dolayısıyla A4 her EV ve her yöntem için sağlanır.
LP en düşük talep için bile olursuz ise talep varış enerjisine eşitlenir.

Her karar yalnızca $t_0$'da mevcut olan bilgiyi kullanır; dolayısıyla bir istasyon işletmecisi işleme hattını çevrim içi olarak çalıştırabilir.
LP olurlu ise, başka bir EV gelmemesi koşuluyla, öngörülen limitler altında nominal plan enerjilerinden başlayarak tüm hedefleri karşılayan bir açık döngü çizelge vardır.
İşleme hattı hiçbir denetleyici için A5'i garanti etmez; çünkü gerçekleşen enerjiler nominal planın altında kalabilir, tahminler gürültülüdür ve sonraki varışlar yük ekler.
Bu nedenle Teorem~\ref{M-thm:invariance} kapsamına girmez; etkisini kabul kontrolü içermeyen değerlendirme geçişleri ölçer.

\subsection{Optimumların doğrusal programları}\label{app:oracle}

Her iki tam bilgili LP de epizodun tüm oturumlarını varış ve ayrılış adımları, varış enerjileri ve kabul edilen gereksinimleriyle; tüm fiyatları ve gerçekleşen $C(t)$ ve $\Crev(t)$ limitlerini birlikte bilir.
Her $i$ oturumu ve $t\in\mathcal T_i=\{t_i^{\sm{arr}},\dots,d_i\}$ adımı için değişkenler $p_i^{+}(t),p_i^{-}(t)\ge0$ ve $E_i(t)$'dir; bunlara bir eksik enerji gevşek değişkeni $z_i\ge0$ ve yumuşak LP'de aşırı yük gevşek değişkenleri $o^{\pm}(t)\ge0$ eşlik eder.
$G_i^{\max}=0.01\,(d_i-t_i^{\sm{arr}}+1)$ kWh yuvarlama payıyla kısıtlar şunlardır:
\begin{equation}\label{eq:lp}
\begin{aligned}
&E_i(t_i^{\sm{arr}})=E_i^{\sm{arr}},\qquad E_i(t+1)=E_i(t)+\Delta t\,\big(p_i^{+}(t)-p_i^{-}(t)\big),\qquad E_i^{\min}-G_i^{\max}\le E_i(t+1)\le E_i^{\sm{cap}},\\
&p_i^{+}(t)\le\Pnom_i,\qquad p_i^{-}(t)\le\Pdis_i,\qquad p_i^{+}(t_i^{\sm{arr}})\le\bar P_i(E_i^{\sm{arr}}),\qquad p_i^{+}(t)\le\varrho_i^{\sm{out}}\big(E_i^{\sm{cap}}-E_i(t)\big)\ (t>t_i^{\sm{arr}}),\\
&E_i(d_i+1)+z_i\ge\Ereq_i-G_i^{\max},\qquad -\Crev(t)-o^{-}(t)\le\textstyle\sum_{i}\big(p_i^{+}(t)-p_i^{-}(t)\big)\le C(t)+o^{+}(t),
\end{aligned}
\end{equation}
burada katı LP $o^{\pm}=0$ olarak sabitler.
Bir çizelgenin kârı $\mathrm{Pr}(p)=\sum_t\Delta t\sum_i\big(c^{\sm{dis}}p_i^{-}-c^{\sm{ch}}p_i^{+}-c^{\sm{thr}}(p_i^{+}+p_i^{-})\big)$'dir; enerji akışı maliyeti $c^{\sm{thr}}$ tek bir duyarlılık çalışması dışında sıfırdır.

\begin{lemma}[Şarj gücü azalması satırları]\label{lem:taperrows}
Her EV, Lemma~\ref{lem:a3} koşullarını sağlasın ve $\varrho_{\kappa,i}=(1-e^{-\kappa_i})/\Delta t$ olsun.
Bu durumda her $E\in[E^{\min},E^{\sm{cap}}]$ için $\bar P_i(E)\le\min\{\Pnom_i,\varrho_{\kappa,i}(E_i^{\sm{cap}}-E)\}$'dir.
Eşitlik, tek adımlık geçiş bandı $\varsigma\in(\varsigma_i^{\sm{tr}}-\delta_i,\varsigma_i^{\sm{tr}})$ dışında sağlanır ve $\sup_E\bar P_i(E)/(E_i^{\sm{cap}}-E)=\varrho_{\kappa,i}$'dir.
\end{lemma}

\begin{proof}
Şarj durumu birimlerinde \Eq{eq:taper} gereğince $\bar P\Delta t/E^{\sm{cap}}=\varphi/E^{\sm{cap}}-\varsigma$'dır ve $\varsigma\ge\varsigma^{\sm{tr}}$ için bu kazanç $\min\{\delta,(1-\varsigma)(1-e^{-\kappa})\}$ olup sağ tarafa eşittir.
$\varsigma+\delta\le\varsigma^{\sm{tr}}$ sabit akım bölgesinde kazanç $\delta$'dır ve sağ taraf ancak ve ancak $\delta\le(1-\varsigma)(1-e^{-\kappa})$ ise $\delta$'ya eşittir.
$1-\varsigma\ge1-\varsigma^{\sm{tr}}=m\delta/\kappa$ olduğundan $\kappa\le m(1-e^{-\kappa})$ olması yeterlidir; $\kappa=mz$ ile bu $z\le1-e^{-mz}$, yani $\kappa\le\kappa^{*}(m)$ demektir.
Geçiş bandında $y=\varsigma^{\sm{tr}}-\varsigma\in(0,\delta)$ olsun; böylece $1-\varphi/E^{\sm{cap}}\ge(1-\varsigma^{\sm{tr}})e^{-m(\delta-y)/(1-\varsigma^{\sm{tr}})}$ olur.
$e^{v}\ge1+v$ ve $m\ge1$ ile bu sınırın sağ tarafı en az $e^{-\kappa}(1-\varsigma^{\sm{tr}}+my)\ge e^{-\kappa}(1-\varsigma)$'dır.
Dolayısıyla kazanç en fazla $(1-\varsigma)(1-e^{-\kappa})$'dır ve zaten en fazla $\delta$'dır; dolayısıyla eşitsizlik her yerde sağlanır ve supremum üst bölgedeki eşitlikten çıkar.
\end{proof}

Kod, kesin değerden \%1 daha gevşek, geçerli bir dış satır olan $\varrho_i^{\sm{out}}=1.01\,\varrho_{\kappa,i}$'yi kullanır.
Kodla birlikte yayımlanan bir denetim betiği, 320 yapılandırmada eşitsizliği $1.7\times10^{-13}$ kW duyarlılığında ve bant dışındaki eşitliği doğrulamaktadır.
Bu satırlarla LP'ler, EV2Gym'in şarj limitini geçiş bandının adım ayrıklaştırması ve \%1'lik pay dışında modeller.
Kalan gevşetmeler yuvarlama payı ve eşzamanlı şarj ile deşarjdır; ikincisi, ana ızgaradaki gibi $\eta=1$, $c^{\sm{dis}}=c^{\sm{ch}}$ ve $c^{\sm{thr}}=0$ olduğunda maliyet ve enerji açısından nötrdür ve dolayısıyla kesindir.
Enerji akışı maliyeti içeren ya da negatif fiyatlarda 0.9 deşarj fiyatı çarpanı kullanan duyarlılık çalışmalarında LP'ler geçerli üst sınırlar olarak kalır.
Bu durumda eşzamanlı şarj ve deşarj kullanabilirler ve bu tür adımların sayısı raporlanır.

\subsection{Üst sınırlar}\label{app:bounds}

Katı LP, HiGHS \cite{huangfu2018parallelizing} ile sözlükbilimsel olarak çözülür.
Aşama 1, $G^{\max}=0$ ile en küçük toplam eksik enerjiyi $z^{*}=\min\sum_iz_i$ hesaplar; aşama 2 ise yuvarlama payı ve $\sum_iz_i\le z^{*}+10^{-6}$ ile $\mathrm{Pr}(p)$'yi en büyükleyerek $J^{\sm{LP,ceil}}$'i verir.
$G^{\max}=0$ ile aynı aşama 2, yalnızca \Sec{sec:results-supp} içinde raporlanan $J^{\sm{LP}}$'yi verir.
Buna karşılık yumuşak LP, \Eq{eq:lp} altında $\mathrm{Pr}(p)-w_E\sum_iz_i-w_P\Delta t\sum_t(o^{+}(t)+o^{-}(t))$ ifadesini tek aşamada en büyükleyerek $J^{\sm{pen,LP}}$'yi verir.

\begin{proposition}[Katı optimum]\label{prop:oracle}
Lemma~\ref{lem:a3} yapılandırması geçerli olsun ve $\pi$, raporlanan güçleri her adımda \Eq{M-eq:transformer} koşulunu sağlayan ve kabul edilen gereksinimlere göre gerçekleşen toplam eksik enerjisi en fazla $z^{*}$ olan herhangi bir denetleyici olsun.
Bu durumda $J^{\pi}\le J^{\sm{LP,ceil}}$'dir.
\end{proposition}

\begin{proof}
$\tilde p_i(t)$ raporlanan güçler ve $E_i^{\sm{nom}}(t)=E_i^{\sm{arr}}+\Delta t\sum_{s<t}\tilde p_i(s)$ yalnızca faturalandırılan güçlerin üreteceği enerji olsun.
EV2Gym tam olarak $\tilde p$'yi faturalandırır; dolayısıyla $J^{\pi}$, $p^{\pm}=\tilde p^{\pm}$ noktasındaki $\mathrm{Pr}$'ye eşittir.
Yuvarlama etkin adım başına 0.01 kWh'ten daha az ekler; dolayısıyla yörünge boyunca $E_i^{\sm{real}}(t)-G_i^{\max}\le E_i^{\sm{nom}}(t)\le E_i^{\sm{real}}(t)\le E_i^{\sm{cap}}$'dir.
Bu nedenle düşürülmüş enerji tabanı ve düşürülmüş ayrılış satırı, $z_i=(\Ereq_i-E_i^{\sm{real}}(d_i+1))^{+}$ ile $E^{\sm{nom}}$ için sağlanır; güç satırları da $\tilde p_i\le\bar P_i(E_i^{\sm{real}})\le\Pnom_i$ ve $\tilde p_i\ge-\Pdis_i$ olduğundan sağlanır.
Ayrıca $E^{\sm{nom}}(t_i^{\sm{arr}})=E_i^{\sm{arr}}$ olduğundan ilk adım satırı sağlanır.
Şarj gücü azalması satırları için EV2Gym hiçbir zaman $\varphi(E)-E$'den fazlasını teslim etmez; dolayısıyla Lemma~\ref{lem:taperrows} $\tilde p_i(t)\le\varrho_{\kappa,i}(E^{\sm{cap}}-E_i^{\sm{real}}(t))\le\varrho_i^{\sm{out}}(E^{\sm{cap}}-E_i^{\sm{nom}}(t))$ sonucunu verir.
Trafo satırları varsayımın kendisidir ve $\sum_iz_i\le z^{*}$'dır; dolayısıyla yörünge aşama 2 için olurludur ve amaç değeri $J^{\pi}$'ye eşittir.
\end{proof}

\begin{proposition}[Yumuşak cezalı optimum]\label{prop:softoracle}
Lemma~\ref{lem:a3} yapılandırması geçerli olsun ve komutlar \Eq{M-eq:limits} ile verilen fiziksel limitlere kırpılsın.
Bir tekrar üzerinde simüle edilen her yörünge için, eksik enerji paylaşılan kabul edilen gereksinime göre ölçüldüğünde, $J^{\sm{pen}}(\pi)\le J^{\sm{pen,LP}}$'dir.
\end{proposition}

\begin{proof}
$p^{\pm}=\tilde p^{\pm}$, $E=E^{\sm{nom}}$, $z_i=\Delta E_i^{\sm{sh}}$ alınsın ve $o^{\pm}(t)$ her iki yöndeki gerçekleşen aşırı yüklere eşit olsun.
Enerji, güç, ilk adım ve şarj gücü azalması satırları Önerme~\ref{prop:oracle} ispatındaki gibi sağlanır; çünkü bunlar denetleyicinin uyduğu herhangi bir kısıtı değil, yalnızca EV2Gym'in fiziğini kullanır.
Ayrılış satırı için $E^{\sm{nom}}(d_i+1)+z_i\ge E^{\sm{real}}(d_i+1)-G_i^{\max}+z_i\ge\Ereq_i-G_i^{\max}$'dir; trafo satırları ise $o^{\pm}$ tanımı gereği sağlanır.
Bu noktada amaç değeri, aynı faturalandırılan güçler, eksik enerji ve aşırı yük enerjisiyle $J^{\sm{pen}}(\pi)$'ye eşittir.
\end{proof}

Sınır denetleyicinin değil ölçümün bir özelliğidir; dolayısıyla eksik enerji ya da aşırı yük içeren denetleyicileri, konformal ve kabul kontrolsüz geçişleri ve harici karşılaştırma yöntemlerini kapsar.
Sarmalayıcıdaki her geçiş fiziksel kırpmayı uygular ve Gurobi eMPC bağdaştırıcısı da komutlarını aynı şekilde kırpar.
Kodla birlikte yayımlanan bir birim testi, rastgele eylemler de dâhil olmak üzere küçük tekrarlarda tüm kural tabanlı yöntemler için $J^{\sm{pen}}\le J^{\sm{pen,LP}}$ koşulunu denetler.
EV2Gym içinde bir denetleyici faturalandırılmayan yuvarlama kazancı yoluyla $J^{\sm{LP}}$'yi aşabilir; bu nedenle ana metin garantinin maliyeti için $J^{\sm{LP,ceil}}$'i kullanır ve $J^{\sm{LP}}$'yi yalnızca burada raporlar.
